\chapter{Metodologia}
\label{cap:metodologia}

\section{Caracterização da pesquisa e etapas}

A pesquisa é de natureza aplicada e tecnológica. O escopo central é o desenvolvimento de um software e a avaliação experimental de seus módulos. O percurso metodológico divide-se nas etapas a seguir.

\begin{enumerate}
	\item \textbf{Levantamento de requisitos}: entrevistas com alunos e professores da LORD-UEPG e análise detalhada da ficha clínica em papel, para compreender o fluxo de atendimento e mapear as variáveis de DTM e prótese (\autoref{sec:requisitos}).
	\item \textbf{Projeto e desenvolvimento do sistema}: definição da arquitetura, da modelagem de dados híbrida e dos mecanismos de segurança, seguida da implementação iterativa das interfaces e serviços (\autoref{sec:arquitetura} a \autoref{sec:seguranca}).
	\item \textbf{Módulo de estruturação por IA}: implementação da transcrição de fala e da estruturação da anamnese com modelos executados localmente, sob revisão humana obrigatória (\autoref{sec:metodo-ia}).
	\item \textbf{Módulo analítico}: cálculo de indicadores de evolução clínica, análises estatísticas e painéis de visualização (\autoref{sec:metodo-analitico}).
	\item \textbf{Avaliação}: experimentos com casos clínicos sintéticos para medir a qualidade da estruturação automática e o efeito da captação por áudio, teste de usabilidade com os usuários finais em ambiente simulado, validação do pipeline analítico e avaliação técnica de desempenho e segurança (\autoref{sec:protocolo}).
\end{enumerate}

\section{Levantamento de requisitos}
\label{sec:requisitos}

Os requisitos foram levantados em reuniões com membros e professores da LORD-UEPG, nas quais foram catalogadas as principais restrições e necessidades do acompanhamento de pacientes com DTM e usuários de prótese. \todo{registrar datas, número de reuniões e participantes; incluir as decisões sobre a ficha de retorno e as categorias de tratamento.}

\subsection{Análise da ficha clínica}
\label{sec:ficha}

A ficha clínica da Liga (Anexo~\ref{anexo:ficha}) foi decomposta em seções, e cada seção foi classificada quanto à natureza da informação. Essa classificação responde a uma questão central do projeto: qual é a forma mais adequada de captar cada tipo de dado. O resultado está no \autoref{qua:ficha-captacao}.

\begin{quadro}[htbp]
	\caption{Seções da ficha clínica, natureza da informação e forma de captação adotada}
	\label{qua:ficha-captacao}
	\footnotesize
	\begin{tabularx}{\textwidth}{|L{4.3cm}|L{3cm}|Y|}
		\hline
		\textbf{Seção da ficha (questões)} & \textbf{Natureza} & \textbf{Captação no sistema} \\ \hline
		Identificação, TCLE e número do prontuário & Cadastral & Formulário digital \\ \hline
		Queixa principal; histórias médica, cirúrgica, de alergias, de medicamentos e psicossocial (1--6); expectativas com a prótese (48--49); observações & Narrativa livre & Ditado pelo acadêmico, com estruturação por IA e revisão humana \\ \hline
		Hábitos, condições bucais e questionário de DTM (7--36); uso de prótese (46--47) & Autorrelato com respostas fechadas e escalas de 0 a 10 & Formulário digital otimizado para \textit{tablet} \\ \hline
		Exame de DTM (37--45); exames extraoral e intraoral (50--75); odontograma; condição periodontal & Exame clínico, medidas e mapas & Componentes visuais interativos (odontograma e mapa facial); ditado opcional durante o exame \\ \hline
		Características relatadas (RDC/TMD); Eixo I -- Diagnóstico; plano de tratamento & Decisão clínica & Exclusivamente humano; bloqueado para a IA \\ \hline
	\end{tabularx}
	\fonte{Os autores (2026).}
\end{quadro}

A captação é, portanto, híbrida. O áudio é empregado onde a informação é naturalmente narrativa ou onde as mãos do acadêmico estão ocupadas e enluvadas durante o exame, situação em que tocar a tela contraria as boas práticas de biossegurança. Formulários são empregados onde as respostas são fechadas e rápidas de marcar. Em todos os casos, o áudio é descartado logo após a transcrição. A comparação objetiva entre as formas de captação é um dos experimentos do trabalho (\autoref{sec:exp-audio}).

\subsection{Requisitos funcionais e não funcionais}

Os principais requisitos estão listados no \autoref{qua:requisitos}.

\begin{quadro}[htbp]
	\caption{Requisitos funcionais (RF) e não funcionais (RNF)}
	\label{qua:requisitos}
	\footnotesize
	\begin{tabularx}{\textwidth}{|L{1.2cm}|Y|}
		\hline
		\textbf{ID} & \textbf{Descrição} \\ \hline
		RF01 & Autenticar usuários e controlar o acesso por papel (aluno, professor, pesquisador, administrador) \\ \hline
		RF02 & Cadastrar e buscar pacientes, registrando os consentimentos (TCLE, gravação de áudio, pesquisa) \\ \hline
		RF03 & Registrar atendimentos de primeira consulta e de retorno, com as seções da ficha clínica \\ \hline
		RF04 & Registrar o odontograma e o mapa de palpação por meio de componentes gráficos \\ \hline
		RF05 & Permitir que o professor revise, devolva com comentários ou aprove as fichas preenchidas pelos alunos \\ \hline
		RF06 & Tornar imutável a ficha aprovada; correções geram nova versão (adendo) \\ \hline
		RF07 & Transcrever o ditado do acadêmico e sugerir o preenchimento dos campos narrativos, exigindo aceite humano campo a campo \\ \hline
		RF08 & Exibir a evolução clínica do paciente em painéis (dor, abertura bucal, palpação, tratamentos) \\ \hline
		RF09 & Exibir painéis de coorte e exportar dados pseudonimizados para pesquisa \\ \hline
		RF10 & Registrar em trilha de auditoria todas as leituras e alterações de prontuários \\ \hline
		RNF01 & Conformidade com a LGPD: minimização, pseudonimização, criptografia e controle de acesso \\ \hline
		RNF02 & Execução integral em infraestrutura da UEPG, sem envio de dados clínicos a serviços externos \\ \hline
		RNF03 & Uso exclusivo de ferramentas gratuitas e de código aberto \\ \hline
		RNF04 & Interface utilizável em \textit{tablet} e computador (\textit{web} responsiva) \\ \hline
		RNF05 & Implantação reprodutível por contêineres \\ \hline
		RNF06 & Evolução da ficha sem perda ou migração forçada dos registros antigos (esquema versionado) \\ \hline
	\end{tabularx}
	\fonte{Os autores (2026).}
\end{quadro}

\subsection{Perfis de usuário e fluxo de atendimento}

O fluxo reproduz a dinâmica pedagógica da Liga, em que os atendimentos são realizados por acadêmicos sob supervisão direta dos professores, conforme estabelece o TCLE. O aluno preenche a ficha, que permanece como \emph{rascunho}. Ao concluir, envia-a para \emph{revisão}. O professor pode devolvê-la com comentários ou \emph{aprová-la}. A ficha aprovada torna-se imutável, e qualquer correção posterior gera uma nova versão vinculada à original, preservando o histórico exigido de um prontuário. O perfil \emph{pesquisador} acessa apenas dados pseudonimizados, e o \emph{administrador} gerencia usuários sem acesso a dados clínicos. A matriz completa de permissões está no Apêndice~\ref{apendice:rbac}.

\begin{figure}[htbp]
	\caption{Diagrama de estados da ficha clínica}
	\label{fig:estados-ficha}
	\centering
	\figtodo{diagrama de estados: Rascunho $\rightarrow$ Em revisão $\rightarrow$ Aprovada; Em revisão $\rightarrow$ Rascunho (devolvida); Aprovada $\rightarrow$ nova versão (adendo)}
	\fonte{Os autores (2026).}
\end{figure}

\section{Materiais e tecnologias}
\label{sec:tecnologias}

Todas as ferramentas empregadas são gratuitas e, com exceção dos serviços de hospedagem de código, de código aberto. O \autoref{qua:stack} resume as escolhas e suas justificativas.

\begin{quadro}[htbp]
	\caption{Tecnologias empregadas}
	\label{qua:stack}
	\footnotesize
	\begin{tabularx}{\textwidth}{|L{2.8cm}|L{4.8cm}|Y|}
		\hline
		\textbf{Camada} & \textbf{Tecnologia} & \textbf{Justificativa} \\ \hline
		Interface & React, TypeScript, Vite, Tailwind CSS, React Router, TanStack Query & Componentização de formulários extensos \cite{parida2024}; formulários gerados a partir do esquema \\ \hline
		Servidor principal & Node.js, Fastify, postgres.js & Entrada e saída não bloqueante para requisições concorrentes \cite{parida2024} \\ \hline
		Banco de dados (evolução) & Migrações em SQL versionadas & Acesso direto a RLS, colunas geradas, \texttt{pgcrypto} e gatilhos \\ \hline
		Banco de dados & PostgreSQL 17, JSONB, \texttt{pgcrypto} & Integridade relacional com flexibilidade documental \cite{postgres-json}; criptografia de campos \\ \hline
		Serviço de IA & Python, FastAPI, faster-whisper, Ollama & Transcrição e estruturação executadas localmente \cite{fasterwhisper,ollama} \\ \hline
		Serviço analítico & Python, Pandas, SciPy & Transformação e análise estatística de dados tabulares \cite{zalke2026} \\ \hline
		Infraestrutura & Docker Compose, Caddy (HTTPS) & Implantação reprodutível no servidor da instituição \\ \hline
		Qualidade & Vitest, Playwright, pytest, OWASP ZAP & Testes automatizados e varredura de segurança \cite{zap} \\ \hline
	\end{tabularx}
	\fonte{Os autores (2026).}
\end{quadro}

\section{Arquitetura do sistema}
\label{sec:arquitetura}

A arquitetura separa a interface, o servidor principal e dois serviços especializados em Python (\autoref{fig:arquitetura}). A interface React comunica-se exclusivamente com a API em Node.js, que concentra a autenticação, o controle de acesso, as regras do prontuário e a auditoria. O serviço de IA recebe da API o áudio ou o texto de uma seção da ficha e devolve uma sugestão estruturada, sem acesso direto ao banco. O serviço analítico conecta-se ao banco com um papel somente de leitura, restrito a um esquema de visões pseudonimizadas.

\begin{figure}[htbp]
	\caption{Arquitetura do sistema}
	\label{fig:arquitetura}
	\centering
	\figtodo{diagrama de contêineres (C4): Navegador/Tablet $\rightarrow$ Caddy (HTTPS) $\rightarrow$ API Node.js $\rightarrow$ PostgreSQL; API $\rightarrow$ Serviço de IA (faster-whisper + Ollama); Serviço analítico $\rightarrow$ esquema pseudonimizado (somente leitura)}
	\fonte{Os autores (2026).}
\end{figure}

\subsection{Infraestrutura e guarda dos dados}

Como o TCLE atribui à UEPG a guarda do prontuário e os dados de saúde são dados sensíveis \cite{lgpd2018}, o sistema foi projetado para ser executado integralmente em infraestrutura institucional. Todos os componentes, inclusive os modelos de IA, são empacotados em contêineres e implantados por um único arquivo Docker Compose em um servidor da universidade. Assim, nenhum dado clínico é enviado a serviços externos. Planos gratuitos de nuvem foram utilizados apenas no desenvolvimento e nos experimentos, e exclusivamente com dados sintéticos. \todo{indicar o servidor efetivamente disponibilizado (DEINFO, LCAD ou NTI) e suas especificações.}

Durante o período deste trabalho, a ficha em papel e o TCLE assinado permanecem como documentos oficiais, e o sistema opera como registro complementar. A eliminação do papel depende de assinatura digital com certificação ICP-Brasil \cite{cfo91-2009}, tratada como trabalho futuro.

\section{Modelagem de dados}
\label{sec:modelagem}

\subsection{Esquema versionado da ficha}

A decisão central da modelagem é descrever a ficha clínica \emph{uma única vez}, como um documento JSON Schema versionado. O mesmo esquema é usado para (i) gerar os formulários da interface, (ii) validar os dados recebidos pela API, (iii) restringir a saída do modelo de linguagem no módulo de IA e (iv) definir a conversão dos documentos em colunas para o módulo analítico. Cada campo carrega metadados próprios, como a indicação \texttt{x-somente-humano}, que impede o preenchimento pela IA. Quando a Liga alterar a ficha, uma nova versão do esquema é criada, e os registros antigos permanecem válidos segundo a versão com que foram criados. Um trecho do esquema está no Apêndice~\ref{apendice:schema}.

\subsection{Modelo híbrido relacional e documental}

Entidades com estrutura estável e relações que exigem integridade referencial, como usuários, pacientes, atendimentos, consentimentos e tratamentos, são modeladas como tabelas relacionais. O conteúdo das seções da ficha, o odontograma e o mapa de palpação são armazenados em colunas \texttt{JSONB}, conforme a justificativa da \autoref{sec:jsonb}. O \autoref{qua:tabelas} descreve as principais tabelas, e a \autoref{fig:der} apresenta o diagrama entidade-relacionamento.

\begin{quadro}[htbp]
	\caption{Principais tabelas do banco de dados}
	\label{qua:tabelas}
	\footnotesize
	\begin{tabularx}{\textwidth}{|L{3cm}|Y|}
		\hline
		\textbf{Tabela} & \textbf{Conteúdo} \\ \hline
		\texttt{usuario} & Nome, e-mail, \textit{hash} da senha (Argon2id), papel, segundo fator de autenticação \\ \hline
		\texttt{paciente} & Número do prontuário, nome, data de nascimento, sexo; CPF, RG, contatos e endereço cifrados \\ \hline
		\texttt{consentimento} & Tipo (TCLE, gravação, pesquisa), versão do termo, data e documento digitalizado \\ \hline
		\texttt{atendimento} & Paciente, data, tipo (primeira consulta ou retorno), aluno, professor e situação \\ \hline
		\texttt{ficha} & Atendimento, seção, versão do esquema, dados (\texttt{JSONB}), situação, autor e aprovador \\ \hline
		\texttt{odontograma} & Atendimento e estado de cada dente e face (\texttt{JSONB}) \\ \hline
		\texttt{tratamento} & Paciente, categoria padronizada (placa oclusal, fisioterapia, farmacológico, prótese etc.), início e fim \\ \hline
		\texttt{sugestao\_ia} & Origem (texto ou áudio), transcrição, modelo, versão do \textit{prompt}, saída (\texttt{JSONB}), latência e decisão humana por campo \\ \hline
		\texttt{auditoria} & Usuário, ação, entidade, diferenças (\texttt{JSONB}), endereço IP e data; apenas inserção \\ \hline
	\end{tabularx}
	\fonte{Os autores (2026).}
\end{quadro}

\begin{figure}[htbp]
	\caption{Diagrama entidade-relacionamento}
	\label{fig:der}
	\centering
	\figtodo{DER completo exportado do banco (dbdiagram.io ou DBeaver)}
	\fonte{Os autores (2026).}
\end{figure}

Os campos mais consultados pelo módulo analítico são expostos como colunas geradas a partir do documento \texttt{JSONB}, indexáveis como colunas comuns, conforme o \autoref{cod:coluna-gerada}. Essa técnica evita duplicação manual de dados e mantém o documento como fonte única.

\begin{lstlisting}[caption={Coluna gerada a partir do documento JSONB}, language=SQL, label=cod:coluna-gerada]
ALTER TABLE ficha
  ADD COLUMN eva_atual smallint
  GENERATED ALWAYS AS ((dados ->> 'q16_eva_atual')::smallint) STORED;

CREATE INDEX idx_ficha_eva ON ficha (eva_atual);
CREATE INDEX idx_ficha_dados ON ficha USING gin (dados jsonb_path_ops);
\end{lstlisting}

A escolha entre o modelo híbrido e a normalização completa é avaliada experimentalmente (\autoref{sec:exp-tecnica}).

O esquema do banco é definido por migrações em SQL, aplicadas em ordem e registradas em uma tabela de controle, e a API o acessa por consultas parametrizadas com a biblioteca postgres.js. Um mapeador objeto-relacional (ORM) foi considerado e descartado. Os requisitos de segurança dependem de recursos específicos do PostgreSQL: políticas de \textit{Row-Level Security}, colunas geradas a partir de JSONB, funções do \texttt{pgcrypto}, gatilhos de imutabilidade e papéis com privilégios distintos. Esses recursos ficam fora do modelo dos ORMs e acabariam escritos em SQL de qualquer forma, com o custo adicional de uma camada intermediária.

\section{Segurança e conformidade com a LGPD}
\label{sec:seguranca}

As medidas de segurança foram definidas a partir do marco legal (\autoref{sec:marco-legal}), das recomendações da literatura \cite{shojaei2024,tertulino2023} e da lista OWASP Top 10 \cite{owasp-top10}. O \autoref{qua:seguranca} relaciona cada requisito à sua implementação.

\begin{quadro}[htbp]
	\caption{Requisitos de segurança e implementação correspondente}
	\label{qua:seguranca}
	\footnotesize
	\begin{tabularx}{\textwidth}{|L{3.2cm}|Y|}
		\hline
		\textbf{Requisito} & \textbf{Implementação} \\ \hline
		Autenticação & Senhas com Argon2id (parâmetros recomendados pela OWASP); sessão em \textit{cookie} \texttt{HttpOnly}, \texttt{Secure} e \texttt{SameSite=Strict}, com o \textit{token} armazenado apenas como \textit{hash}; bloqueio por 15 minutos após cinco tentativas; mesma resposta para e-mail inexistente e senha incorreta. Segundo fator (TOTP) para professores previsto \\ \hline
		Controle de acesso & Controle por papéis (RBAC) na API e \textit{Row-Level Security} no PostgreSQL. A API conecta com um papel que não é dono das tabelas e informa o usuário a cada transação. O aluno só enxerga os pacientes que cadastrou ou atendeu, o professor enxerga todos, e administrador e pesquisador não enxergam dados clínicos identificados \\ \hline
		Confidencialidade em trânsito & HTTPS em todas as comunicações \\ \hline
		Confidencialidade em repouso & Campos identificadores cifrados com \texttt{pgcrypto}; disco do servidor cifrado; segredos fora do repositório de código \\ \hline
		Pseudonimização (LGPD, art.~13) & Esquema analítico com visões que substituem o identificador do paciente por um código HMAC com chave secreta mantida em esquema inacessível aos papéis de aplicação. As visões usam faixas etárias em vez de datas de nascimento, dias desde a primeira consulta em vez de datas, e omitem nome, documentos, contatos e todos os campos de texto livre \\ \hline
		Integridade do prontuário & Gatilhos no banco impedem a alteração de fichas aprovadas e a exclusão de qualquer ficha; correções são registradas como adendos versionados \\ \hline
		Anonimização em exportações & Exportação somente a partir do esquema analítico, com verificação de k-anonimato (k $\geq$ 5) nas combinações de atributos quase identificadores \\ \hline
		Rastreabilidade & Trilha de auditoria de leituras e alterações, sem permissão de alteração ou exclusão para a aplicação \\ \hline
		Minimização & Descarte do áudio após a transcrição; coleta restrita aos campos da ficha \\ \hline
		Disponibilidade & Cópias de segurança diárias cifradas em armazenamento institucional, com teste periódico de restauração \\ \hline
	\end{tabularx}
	\fonte{Os autores (2026).}
\end{quadro}

\section{Módulo de estruturação da anamnese por IA}
\label{sec:metodo-ia}

O módulo transforma o relato do acadêmico em sugestões de preenchimento da ficha, em quatro passos (\autoref{fig:fluxo-ia}):

\begin{enumerate}
	\item \textbf{Transcrição}: o áudio é transcrito localmente pelo faster-whisper \cite{fasterwhisper}, com um \textit{prompt} inicial contendo o vocabulário do domínio (ATM, masseter, crepitação, bruxismo, prótese parcial removível etc.); em seguida, o áudio é descartado;
	\item \textbf{Estruturação}: a transcrição, ou o texto digitado, é enviada a um modelo de linguagem de pesos abertos executado pelo Ollama \cite{ollama}, com temperatura zero e saída restrita ao JSON Schema da seção correspondente, do qual foram removidos os campos marcados como exclusivamente humanos;
	\item \textbf{Validação}: a saída é validada contra o esquema e passa por verificações determinísticas, descritas adiante, que descartam valores sem evidência, com evidência ausente do relato ou incoerentes com a própria evidência;
	\item \textbf{Revisão humana}: a interface exibe cada sugestão destacada, acompanhada do trecho que a justifica, com as opções aceitar, editar e rejeitar. Nenhum valor é gravado sem uma dessas ações.
\end{enumerate}

\begin{figure}[htbp]
	\caption{Fluxo do módulo de estruturação da anamnese}
	\label{fig:fluxo-ia}
	\centering
	\figtodo{diagrama de sequência: Aluno $\rightarrow$ Interface $\rightarrow$ API $\rightarrow$ Serviço de IA (Whisper $\rightarrow$ LLM $\rightarrow$ validação) $\rightarrow$ API $\rightarrow$ tela de revisão $\rightarrow$ ficha (rascunho)}
	\fonte{Os autores (2026).}
\end{figure}

Para cada campo, o modelo devolve o valor e o trecho de evidência, como no \autoref{cod:saida-ia}. A extração é feita por seção da ficha, e não da ficha inteira de uma vez, o que reduz o volume de texto gerado e, consequentemente, o tempo de resposta em servidores sem GPU. As duas estratégias são comparadas no Experimento~A.

\begin{lstlisting}[caption={Exemplo de sugestão gerada para a seção de história médica}, label=cod:saida-ia]
{
  "q3_alergias": {
    "valor": "Dipirona (manchas na pele).",
    "evidencia": "tem alergia a dipirona, deu mancha na pele"
  },
  "q4_usa_medicamento": {
    "valor": true,
    "evidencia": "toma losartana faz uns dois anos"
  },
  "q2_historico_cirurgico": { "valor": null, "evidencia": null }
}
\end{lstlisting}

O formato de saída é gerado automaticamente a partir do esquema da seção: cada campo vira um par \texttt{valor}/\texttt{evidencia}, e o Ollama restringe a decodificação a esse formato (\textit{structured outputs}). A instrução enviada ao modelo é composta por regras fixas e por um guia dos campos, também derivado do esquema, com o título, a orientação de preenchimento e os códigos aceitos de cada campo, acompanhados de seus rótulos. Assim, uma alteração na ficha propaga-se ao \textit{prompt} sem edição manual.

Antes de chegar à tela de revisão, cada valor passa por três verificações determinísticas, implementadas no pacote compartilhado da ficha e cobertas por testes automatizados:

\begin{enumerate}
	\item \textbf{existência de evidência}: valores sem trecho de evidência são descartados; o texto ``null'', que os modelos às vezes escrevem no lugar do valor nulo, é tratado como ausência;
	\item \textbf{evidência presente no relato}: o trecho citado deve constar do relato após normalização de caixa, pontuação e espaços; para tolerar pequenas variações do modelo, admite-se também que ao menos 90\% das palavras significativas do trecho estejam no relato. Aspas e rótulos em volta da citação (``Relato do acadêmico: \ldots'') são removidos antes da comparação. Uma paráfrase, porém, é recusada: ``nega alergias'' não é aceita como evidência quando o relato diz ``não tem alergia a nada'', pois o que não é literal não pode ser conferido. Isso impede que o modelo ``cite'' algo que não foi dito;
	\item \textbf{coerência entre valor e evidência}: em campos de texto, alguma palavra significativa do valor deve aparecer no trecho citado, desconsideradas as flexões de gênero e número (``respiratória'' e ``respiratório''), e um valor de negação (``Nega.'') exige uma negação na evidência. Isso impede, por exemplo, registrar ``Nega.'' em alergias com a evidência ``alergia a dipirona''.
\end{enumerate}

Essas verificações foram refinadas da mesma forma que o \textit{prompt}: o roteiro do Experimento~A grava a saída crua do modelo e a lista de valores descartados com o motivo, o que permitiu identificar descartes indevidos, como o de citações corretas embrulhadas em rótulos, e reavaliar as salvaguardas sem executar o modelo novamente.

Os valores descartados não são gravados na ficha, mas são exibidos ao acadêmico, com o motivo, em uma lista recolhida da tela de revisão. As verificações não substituem a revisão humana: um valor apoiado em evidência verdadeira, porém mal classificado, continua a depender do acadêmico para ser rejeitado.

O \textit{prompt} é versionado, e a versão fica registrada em cada sugestão. Os ajustes do \textit{prompt} foram feitos exclusivamente com relatos de desenvolvimento escritos pelos autores para casos do conjunto de ajuste (\autoref{sec:casos}), nunca com casos de teste. O \autoref{qd:versoes-prompt} resume as versões e o problema observado que motivou cada uma.

\begin{quadro}[htbp]
	\caption{Versões do \textit{prompt} de estruturação}
	\label{qd:versoes-prompt}
	\centering
	\small
	\begin{tabularx}{\textwidth}{|L{1.2cm}|Y|Y|}
		\hline
		\textbf{Versão} & \textbf{Problema observado} & \textbf{Alteração} \\ \hline
		v1 & --- & Regras gerais: não diagnosticar, não deduzir, evidência literal, negações explícitas, somente códigos listados, ignorar comentários irrelevantes \\ \hline
		v2 & O modelo marcou ``respiratória'' a partir do uso de losartana e repetiu a mesma informação em dois campos & Proibição de deduzir doenças a partir de medicamentos; cada informação em um único campo; exemplo resolvido com outro paciente \\ \hline
		v3 & O modelo copiou o título do campo e o rótulo ``Valor:'' do guia para dentro do valor & Guia reescrito com rótulos que não se confundem com conteúdo; regra explícita de que o valor contém apenas a informação do paciente \\ \hline
		v4 & O modelo citou o título do campo como evidência (``Tipo de escova'') ou embrulhou a citação (``Relato do acadêmico: `...' '') & Regra da evidência passa a proibir aspas, prefixos e o uso do título ou da pergunta do campo \\ \hline
	\end{tabularx}
	\fonte{Os autores (2026).}
\end{quadro}

As salvaguardas que asseguram a ausência de finalidade diagnóstica são implementadas no código, e não apenas declaradas:

\begin{itemize}
	\item as seções de diagnóstico (Eixo~I do RDC/TMD), características relatadas e plano de tratamento são marcadas como \texttt{x-somente-humano}, não são enviadas ao modelo e têm escrita rejeitada pela API quando a origem é a IA;
	\item a instrução ao modelo proíbe inferir, diagnosticar ou completar informações ausentes, exigindo \texttt{null} para o que não foi dito, e a API recusa pedidos de sugestão para seções exclusivamente humanas e para fichas fora do estado de rascunho;
	\item escores derivados, como a GCPS, são calculados por fórmula determinística, sem IA;
	\item a tabela \texttt{sugestao\_ia} registra o modelo, a versão do \textit{prompt}, a saída bruta e a decisão humana sobre cada campo (aceita, aceita e editada ou rejeitada), permitindo auditoria e medição da taxa de aceitação; a trilha de auditoria registra apenas contagens, sem o conteúdo do relato.
\end{itemize}

O modelo de referência do desenvolvimento é o Qwen2.5 de 3 bilhões de parâmetros (\texttt{qwen2.5:3b}, quantização Q4\_K\_M) executado pelo Ollama 0.34.3, com temperatura zero, semente fixa e janela de contexto de 8.192 \textit{tokens}. A transcrição usa o modelo \texttt{small} do faster-whisper 1.2.1 em CPU, com quantização \texttt{int8}. \todo{definir o segundo modelo avaliado no Experimento~A (por exemplo, \texttt{qwen2.5:7b} ou \texttt{llama3.1:8b}) e registrar a versão exata.}

\section{Módulo analítico}
\label{sec:metodo-analitico}

O módulo analítico lê exclusivamente o esquema pseudonimizado, converte os documentos das fichas em tabelas com o Pandas e disponibiliza os resultados à interface por meio de uma API. Os indicadores estão organizados em três níveis.

\subsection{Indicadores individuais}

Para cada paciente, ao longo das consultas, são calculados:

\begin{itemize}
	\item intensidade da dor atual, da pior dor e da dor média (questões 16, 17 e 18);
	\item abertura bucal sem dor, máxima sem auxílio e máxima com auxílio, em milímetros (questão 39);
	\item número de sítios dolorosos à palpação e soma das intensidades de 0 a 3 (questões 44 e 45);
	\item número de atividades limitadas pela dor (questão 36);
	\item grau da Escala Graduada de Dor Crônica \cite{vonkorff1992}, conforme as Equações~(\ref{eq:cpi}) e~(\ref{eq:incap}).
\end{itemize}

A intensidade característica da dor (ICD) e o escore de incapacidade (EI) são dados por

\begin{equation}
	\label{eq:cpi}
	\mathrm{ICD} = \frac{Q_{16} + Q_{17} + Q_{18}}{3} \times 10
\end{equation}

\begin{equation}
	\label{eq:incap}
	\mathrm{EI} = \frac{Q_{20} + Q_{21} + Q_{22}}{3} \times 10
\end{equation}

\noindent em que $Q_i$ é a resposta da questão $i$. O EI e o número de dias de afastamento ($Q_{19}$) são convertidos em pontos de incapacidade (0 a 3 cada): de 0 a 6 dias, 0 ponto; de 7 a 14, 1; de 15 a 30, 2; 31 ou mais, 3; EI de 0 a 29, 0 ponto; de 30 a 49, 1; de 50 a 69, 2; 70 ou mais, 3. O grau é 0 na ausência de dor; I se ICD~<~50 e menos de 3 pontos; II se ICD~$\geq$~50 e menos de 3 pontos; III com 3 ou 4 pontos; e IV com 5 ou 6 pontos. \todo{conferir os pontos de corte no manual do RDC/TMD \cite{dworkin1992,pereira2004}.}

\subsection{Indicadores de coorte}

\begin{itemize}
	\item prevalência das características relatadas (estalido, crepitação, apertamento diurno e noturno, rigidez matinal, zumbido) e de cefaleia;
	\item distribuição dos diagnósticos do Eixo~I, registrados pelos professores;
	\item variação da dor entre a consulta inicial e a mais recente, por categoria de tratamento, com teste de Wilcoxon pareado dentro de cada grupo e teste de Kruskal-Wallis entre grupos, acompanhados de tamanho de efeito;
	\item proporção de pacientes com redução da dor de pelo menos 30\% (melhora moderadamente importante) e de pelo menos 50\% (melhora substancial), conforme o IMMPACT \cite{dworkin2008};
	\item correlação de Spearman entre variáveis psicossociais (sono, ansiedade) e intensidade da dor;
	\item taxa de abandono do tratamento, definida como a ausência de retorno por período superior ao combinado.
\end{itemize}

\subsection{Indicadores de qualidade dos dados}

\begin{itemize}
	\item \textbf{completude}: proporção de campos obrigatórios preenchidos por ficha;
	\item \textbf{tempo de preenchimento} por seção;
	\item \textbf{taxa de aceitação} das sugestões da IA, por campo e por seção.
\end{itemize}

\subsection{Visualizações}

São disponibilizados três painéis: (i) do paciente, com a evolução da dor, da abertura bucal e dos tratamentos em linha do tempo e um mapa facial de calor da palpação por consulta; (ii) da coorte, com diagramas de caixa da variação da dor por tratamento e gráficos de prevalência; e (iii) de qualidade dos dados.

A análise da eficácia dos tratamentos exige que as mesmas medidas sejam repetidas nos retornos. Como a ficha atual é de primeira consulta, foi definida com a Liga uma \emph{ficha de retorno} reduzida, com as questões 16 a 22, 39, 44 e 45, a adesão ao tratamento e as intercorrências, além de uma lista fechada de categorias de tratamento. \todo{confirmar com a LORD o conteúdo da ficha de retorno e a lista de tratamentos.}

\section{Protocolo de avaliação}
\label{sec:protocolo}

Para não depender de dados reais de pacientes durante o desenvolvimento, todos os experimentos com IA e com o módulo analítico utilizam dados sintéticos. Dados reais só serão tratados no servidor institucional e mediante as autorizações descritas na \autoref{sec:etica}.

\subsection{Base de casos clínicos sintéticos}
\label{sec:casos}

Foram elaborados 40 casos clínicos fictícios. Cada caso é composto por um \emph{gabarito}, isto é, a ficha inicial completa em formato JSON, e por uma \emph{narrativa} de consulta que contém as informações do gabarito.

Os gabaritos são produzidos por um gerador em Python, com semente aleatória fixa (2026) para garantir a reprodutibilidade. Os casos são distribuídos igualmente entre cinco perfis clínicos: DTM predominantemente muscular, DTM articular, usuário de prótese, DTM associada ao uso de prótese e paciente sem DTM. As respostas são sorteadas sob regras de coerência clínica. Por exemplo, a pior dor nunca é menor que a dor atual, a abertura com auxílio supera a abertura sem auxílio, e o diagnóstico registrado é compatível com os achados de palpação. Todo gabarito é validado contra o esquema da ficha, e as regras de coerência são verificadas por testes automatizados.

Para cada caso, o gerador produz também um roteiro destinado ao membro da Liga que redigirá a narrativa. O roteiro lista os fatos que devem aparecer, os campos que \emph{não} podem ser mencionados, por estarem vazios no gabarito, e de um a três desafios de linguagem sorteados. Os desafios são negações (``nega alergias''), expressões temporais imprecisas (``faz uns dois anos''), autocorreções (``não, é do lado esquerdo''), termos coloquiais (``trava'', ``dentadura''), informações irrelevantes e ordem diferente da ficha. A lista de campos a omitir torna mensurável a taxa de alucinação: qualquer valor extraído para esses campos é, por construção, indevido.

Os casos foram divididos, por sorteio, em um conjunto de ajuste (10 casos), usado para refinar os \textit{prompts}, e um conjunto de teste (30 casos), usado apenas na medição final.

\subsection{Experimento A: qualidade da estruturação automática}
\label{sec:exp-ia}

Os casos de teste são processados por \todo{dois ou três} modelos de linguagem, cada um em duas estratégias (ficha inteira e extração por seção), restritas às seções que admitem ditado: queixa principal, história médica, história odontológica, prótese e observações. Cada campo extraído é comparado ao gabarito e classificado como verdadeiro positivo ($VP$, valor correto), falso positivo ($FP$, valor em campo vazio no gabarito), falso negativo ($FN$, campo com valor deixado vazio), verdadeiro negativo ou erro de valor ($EV$, campo preenchido com valor diferente do gabarito). Como é usual em extração de informação, o erro de valor conta simultaneamente como falso positivo e falso negativo:

\begin{equation}
	P = \frac{VP}{VP+FP+EV}, \qquad R = \frac{VP}{VP+FN+EV}, \qquad F_1 = \frac{2PR}{P+R}.
\end{equation}

Campos categóricos, booleanos e numéricos são comparados por igualdade exata; listas, como conjuntos. Campos de texto livre recebem uma medida automática aproximada, o $F_1$ entre os conjuntos de palavras do valor extraído e do gabarito, sendo considerados corretos a partir de 0,5; a medida definitiva, porém, é a avaliação por dois avaliadores em escala de 0 (inadequado) a 2 (adequado), a partir de uma planilha gerada automaticamente com os pares gabarito--valor extraído. A medida automática penaliza paráfrases corretas (``Não bebe, não fuma e não usa drogas'' diante de ``Nega álcool, drogas e cigarro''), razão pela qual o $F_1$ é informado também separadamente para os campos fechados, de comparação exata, e para os campos de texto. Todas as métricas são calculadas em dois momentos: sobre a saída \emph{bruta} do modelo e sobre a saída \emph{após as salvaguardas}, o que permite quantificar o efeito das verificações determinísticas. O experimento é executado por um \textit{script} que grava cada execução assim que termina, de modo que uma interrupção não perde as execuções concluídas. Além disso, medem-se:

\begin{itemize}
	\item \textbf{taxa de alucinação}: proporção de campos preenchidos pelo modelo cujo valor no gabarito é nulo;
	\item \textbf{taxa de omissão}: proporção de campos com valor no gabarito deixados nulos pelo modelo;
	\item \textbf{validade estrutural}: proporção de respostas que respeitam o esquema;
	\item \textbf{latência} por caso, em CPU e em GPU;
	\item \textbf{reprodutibilidade}: concordância entre três execuções do mesmo caso.
\end{itemize}

\subsection{Experimento B: captação por áudio versus texto}
\label{sec:exp-audio}

\todo{15 a 20} narrativas do conjunto de teste são lidas e gravadas por voluntários em duas condições acústicas: ambiente silencioso e ambiente com ruído de clínica (ruído gravado na própria clínica e sobreposto ao áudio). As gravações são transcritas por modelos Whisper de diferentes tamanhos. A qualidade da transcrição é medida pela taxa de erro de palavras

\begin{equation}
	\mathrm{WER} = \frac{S + D + I}{N},
\end{equation}

\noindent em que $S$, $D$ e $I$ são, respectivamente, as substituições, remoções e inserções necessárias para transformar a transcrição no texto de referência e $N$ é o número de palavras da referência. A WER é calculada também apenas sobre os termos clínicos. Em seguida, as transcrições passam pelo Experimento~A, e a diferença de $F_1$ entre a extração a partir da transcrição e a partir do texto original quantifica a perda introduzida pelo áudio. Registra-se ainda o tempo de fala em comparação com o tempo de digitação do mesmo conteúdo.

\subsection{Experimento C: usabilidade e eficiência}
\label{sec:exp-usabilidade}

\begin{itemize}
	\item \textbf{Participantes}: \todo{8 a 12} acadêmicos da LORD-UEPG e \todo{2 ou 3} professores, estes para o fluxo de revisão. Critério de inclusão: ser membro ativo e já ter preenchido a ficha em papel. Critério de exclusão: ter participado do desenvolvimento do sistema. O tamanho da amostra baseia-se na evidência de que poucos usuários revelam a maior parte dos problemas de usabilidade \cite{nielsen1993}.
	\item \textbf{Desenho}: intrassujeitos. Cada participante registra três casos simulados, apresentados por roteiro gravado ou por um colega no papel de paciente, em três condições: ficha em papel, formulário digital e formulário digital com ditado e IA. A ordem das condições é contrabalanceada por quadrado latino, e os casos têm complexidade equivalente.
	\item \textbf{Instrumentos}: cronometragem por ficha; comparação com o gabarito para medir completude e erros; questionário SUS \cite{brooke1996}; NASA-TLX \cite{hart1988}; entrevista semiestruturada curta (Apêndice~\ref{apendice:usabilidade}).
	\item \textbf{Indicadores}: tempo de preenchimento, completude, número de erros, escore SUS, carga de trabalho, taxa de aceitação das sugestões da IA e problemas de usabilidade classificados por severidade.
	\item \textbf{Análise}: teste de Friedman com comparações de Wilcoxon a posteriori, adequados a amostras pequenas; escore SUS médio com intervalo de confiança, interpretado segundo \citeonline{bangor2008}; categorização das respostas abertas.
\end{itemize}

\subsection{Experimento D: validação do pipeline analítico}
\label{sec:exp-analitico}

Foi gerada uma coorte sintética de \todo{150} pacientes com \todo{4 a 6} consultas cada, com trajetórias de dor \emph{programadas}: por exemplo, redução média de 40\% no grupo com placa oclusal e de 10\% no grupo apenas com orientações, acrescidas de ruído aleatório e de abandono. O módulo analítico é aplicado a essa coorte e deve recuperar os efeitos programados. Por ter efeito conhecido, a coorte valida objetivamente o cálculo dos indicadores e dos testes estatísticos, sem pretender produzir conclusões clínicas.

\subsection{Avaliação técnica}
\label{sec:exp-tecnica}

\begin{itemize}
	\item \textbf{Desempenho da persistência}: o odontograma e o mapa de palpação de \todo{1.000 e 10.000} pacientes sintéticos são armazenados no modelo híbrido (\texttt{JSONB}) e em um modelo totalmente normalizado; comparam-se o tempo de leitura do prontuário completo (\texttt{EXPLAIN ANALYZE}), o tempo de escrita e o espaço em disco;
	\item \textbf{Segurança}: varredura com o OWASP ZAP \cite{zap} antes e depois das correções e testes automatizados que verificam que cada papel não acessa recursos não autorizados;
	\item \textbf{Qualidade do código}: cobertura dos testes unitários e de ponta a ponta executados em integração contínua.
\end{itemize}

\section{Aspectos éticos}
\label{sec:etica}

O Experimento~C envolve participantes humanos (acadêmicos e professores) e, por isso, \todo{foi submetido ao Comitê de Ética em Pesquisa da UEPG por meio da Plataforma Brasil, CAAE n.º ...}, observadas as Resoluções CNS n.º 466/2012 e 510/2016 \cite{cns466-2012,cns510-2016}. Os participantes assinam termo de consentimento próprio e suas respostas são analisadas de forma anônima. Nenhum dado real de paciente é utilizado nos experimentos A, B e D. O uso futuro do ditado em atendimentos reais exige termo de consentimento adicional para gravação de áudio, uma vez que o TCLE vigente não contempla essa finalidade.

\section{Cronograma}

As atividades previstas estão na \autoref{tb:cronograma}.

\begin{table}[htbp]
	\caption{Cronograma de atividades (2026)}
	\label{tb:cronograma}
	\footnotesize
	\begin{tabularx}{\textwidth}{|L{3.2cm}|Y|}
		\hline
		\textbf{Período} & \textbf{Atividades} \\ \hline
		28/09 a 04/10 & Reunião com a LORD; submissão ao CEP; ambiente de desenvolvimento; esquema JSON da ficha (versão 1) \\ \hline
		05/10 a 18/10 & API (autenticação, papéis, pacientes, atendimentos, fichas, auditoria); formulários gerados pelo esquema; odontograma e mapa de palpação; fluxo de revisão \\ \hline
		19/10 a 01/11 & Serviço de IA e tela de revisão das sugestões; base de casos sintéticos; Experimento A \\ \hline
		02/11 a 08/11 & Gravações e Experimento B; serviço analítico, coorte sintética e painéis (Experimento D) \\ \hline
		09/11 a 15/11 & Experimento C; avaliação técnica; implantação no servidor da UEPG \\ \hline
		16/11 a 29/11 & Redação dos resultados e da discussão; revisão geral e normalização \\ \hline
		30/11 a 06/12 & Correções finais e preparação da defesa \\ \hline
	\end{tabularx}
	\fonte{Os autores (2026).}
\end{table}

\todo{ajustar o cronograma à data da defesa.}
